% !TEX program = xelatex
\documentclass[10pt,a4paper]{article}

% ---- Пакеты для XeLaTeX и кириллицы ----
\usepackage{fontspec}
\usepackage{polyglossia}
\setmainlanguage{russian}
\setotherlanguage{english}

% ---- Шрифт с кириллицей ----
\setmainfont{Liberation Serif}
\setsansfont{Liberation Sans}
\setmonofont{Liberation Mono}

% ---- Поля и геометрия ----
\usepackage[top=1.5cm, bottom=1.5cm, left=1.5cm, right=1.5cm]{geometry}
\usepackage{parskip}
\setlength{\parindent}{0pt}

% ---- Продвинутые ссылки ----
\usepackage{hyperref}
\hypersetup{
    colorlinks=true,
    linkcolor=blue,
    urlcolor=blue,
    citecolor=blue
}

% ---- Полезные пакеты ----
\usepackage{enumitem}
\setlist{nosep, left=0pt}
\usepackage{tabularx}
\usepackage{array}
\usepackage{booktabs}
\usepackage{xcolor}
\usepackage{titlesec}

% ---- Настройка заголовков разделов ----
\titleformat{\section}{\Large\bfseries\uppercase}{}{0em}{}[\vspace{-0.3em}\rule{\textwidth}{0.5pt}]
\titlespacing*{\section}{0pt}{0.8em}{0.3em}

\titleformat{\subsection}{\large\bfseries}{}{0em}{}
\titlespacing*{\subsection}{0pt}{0.4em}{0.2em}

% ---- Пользовательские команды ----
\newcommand{\projectEntry}[3]{\textbf{#1} \hfill \textit{#2}\par #3\par\vspace{0.3em}}

% ---- Общие данные (имя, телефон, email и т.д.) ----
% ============================================================
% ОБЩИЕ ПЕРСОНАЛЬНЫЕ ДАННЫЕ
% ============================================================
% Меняешь здесь — меняется во всех резюме

\newcommand{\PersonName}{Кузык Иван Романович}
\newcommand{\PersonPhone}{+7 (916) 574-52-99}
\newcommand{\PersonEmail}{vanya.kuzyk1@gmail.com}
\newcommand{\PersonGitHub}{github.com/vanya890}
\newcommand{\PersonCity}{г.~Москва, м.~Бульвар Дмитрия Донского}
\newcommand{\PersonCitizenship}{РФ}
\newcommand{\PersonDrivingLicense}{Категория B}

% ============================================================
% КОМАНДА ДЛЯ ЗАГОЛОВКА РЕЗЮМЕ
% ============================================================
% Использование: \makeHeader{Должность}
\newcommand{\makeHeader}[1]{%
\begin{center}
    {\Huge\bfseries \PersonName}\\[0.3em]
    {\large #1}\\[0.5em]
    \begin{tabular}{ll}
        Телефон: & \PersonPhone \\
        Email:   & \href{mailto:\PersonEmail}{\PersonEmail} \\
        GitHub:  & \href{https://\PersonGitHub}{\PersonGitHub} \\
        Город:   & \PersonCity \\
        Гражданство: & \PersonCitizenship
    \end{tabular}
\end{center}
}

% ============================================================
% ОБЩЕЕ ОБРАЗОВАНИЕ
% ============================================================
% Использование: \makeEducation
\newcommand{\makeEducation}{%
\section{Образование}

\projectEntry{Московский политехнический университет --- Факультет информационных технологий}{2022--2026}
{Бакалавриат, направление «Информатика и вычислительная техника», специализация «Киберфизические системы». ВКР «Разработка образовательного 3D-симулятора робототехники» --- отлично. Средний балл --- 5.0.}

\projectEntry{ГБПОУ «Образовательный комплекс «Юго-Запад»}{2018--2022}
{Среднее специальное образование, специальность «Информационные системы и программирование». Диплом с отличием. Демонстрационный экзамен по компетенции «Разработка компьютерных игр и мультимедийных приложений» --- отлично.}
}

% ============================================================
% ОБЩАЯ ДОПОЛНИТЕЛЬНАЯ ИНФОРМАЦИЯ
% ============================================================
% Использование: \makeFooter{Ключевые компетенции}{Личные качества}
\newcommand{\makeFooter}[2]{%
\section{Дополнительная информация}

\begin{itemize}
    \item \textbf{GitHub:} \href{https://\PersonGitHub}{\PersonGitHub} --- #1
    \item \textbf{Ключевые компетенции:} #2
    \item \textbf{Личные качества:} Обучаемость, ответственность, умение работать с большим объёмом информации, системное мышление, ориентация на результат
    \item \textbf{Водительские права:} \PersonDrivingLicense
\end{itemize}
}

\begin{document}

% ============================================================
% ЗАГОЛОВОК
% ============================================================
\makeHeader{Специалист по AI/ML / Промпт-инженер}

% ============================================================
% О СЕБЕ
% ============================================================
\section{О себе}

Специалист в области искусственного интеллекта и машинного обучения. Имею практический опыт работы с различными AI-сервисами и моделями, включая Claude, Yandex AI, Google AI Studio, Groq, Qwen Coder. Занимался развёртыванием локальных LLM (Unsloth, LM Studio, Ollama), обучением LoRA-моделей, генерацией изображений (Stable Diffusion, ComfyUI). Разрабатывал собственные нейросетевые модели (перцептроны, CNN) и AI-агентов для автоматизации. Владею промпт-инжинирингом, составлением JSON schema для структурированного вывода моделей.

% ============================================================
% ОПЫТ РАБОТЫ (ПРОЕКТНАЯ ДЕЯТЕЛЬНОСТЬ)
% ============================================================
\section{Опыт работы}

\subsection{Проектная деятельность / Фриланс}

\projectEntry{Разработка AI-агента для автоматической публикации статей}{2025}
{Создал бота для сайта, который автоматически генерировал и публиковал статьи на актуальные темы. Использовал различные языковые модели: \textbf{Claude, Yandex AI, Google AI Studio, Groq, Qwen Coder}. Настроил пайплайн: генерация контента $\to$ проверка $\to$ автоматическая публикация.}

\projectEntry{Fullstack-разработчик (AI Chat Game)}{2024–2025}
{Создал веб-приложение для чата с AI-персонажами на \textbf{React 18 + TypeScript + Vite}. Интегрировал поддержку нескольких AI-провайдеров: \textbf{OpenAI, Ollama, LM Studio, Groq}. Реализовал систему памяти (кратко- и долгосрочную), дашборд аналитики токенов с разбивкой по моделям.}

\projectEntry{Fullstack-разработчик (InfinityAI Novel)}{2025–2026}
{Разработал AI-визуальную новеллу на \textbf{React 19 + TypeScript + Vite 7 + Tailwind CSS 4}. Интегрировал генерацию сюжета через \textbf{OpenRouter API} и локальные модели \textbf{LM Studio} с \textbf{JSON Schema structured output}. Реализовал генерацию портретов персонажей через \textbf{ComfyUI}.}

\projectEntry{Исследователь (AI/ML — нейросети и обучение)}{2024–2026}
{Выполнил цикл лабораторных и учебных проектов по искусственному интеллекту и машинному обучению:
\begin{itemize}
    \item Реализовал самописный \textbf{перцептрон} для распознавания чисел на Python (NumPy)
    \item Разработал систему определения входящих/выходящих людей из автобуса с использованием \textbf{OpenCV} и компьютерного зрения
    \item Реализовал управление беспилотным 2D-роботом в гоночной симуляции (обучение с подкреплением)
    \item Работал с \textbf{Scikit-learn, NumPy, Matplotlib} для задач классификации и визуализации
\end{itemize}}

\projectEntry{Работа с моделями генерации изображений}{2024–2025}
{Настраивал пайплайны генерации изображений в \textbf{Stable Diffusion} и \textbf{ComfyUI}. Работал с \textbf{ControlNet} для управления генерацией. Обучал собственные \textbf{LoRA-модели} на подготовленных датасетах.}

\projectEntry{Развёртывание локальных LLM}{2024–2025}
{Устанавливал и настраивал локальные языковые модели с использованием \textbf{Unsloth, LM Studio, Ollama, llama.cpp}. Настраивал инференс, оптимизировал производительность, работал с открытыми весами моделей.}

\projectEntry{SEO-оптимизация с использованием AI (Stroy-Materiali-24.ru)}{2024}
{Разработал инструменты для SEO-оптимизации сайта строительных материалов. Использовал \textbf{AI-промпты} для генерации контента, \textbf{Яндекс API} и \textbf{Яндекс.Вебмастер} для анализа поисковых запросов.}

% ============================================================
% ПРОМПТ-ИНЖИНИРИНГ
% ============================================================
\section{Промпт-инжиниринг}

\begin{itemize}
    \item Разработка структурированных промптов для разных LLM (Claude, GPT, Yandex AI, Qwen)
    \item Создание \textbf{JSON schema} для управления структурированным выводом моделей
    \item Оптимизация промптов под конкретные задачи: генерация кода, контента, анализ данных
    \item Работа с system prompt, few-shot learning, chain-of-thought
    \item Интеграция промптов в приложения через API
\end{itemize}

% ============================================================
% ОБРАЗОВАНИЕ
% ============================================================
\makeEducation

% ============================================================
% НАВЫКИ
% ============================================================
\section{Навыки}

\begin{tabular}{p{3.5cm} p{12cm}}
    \textbf{AI-сервисы:} & Claude, Cursor, Cline, Antigravity, Google AI Studio, Yandex AI, Grok, Qwen Coder, OpenAI API \\
    \textbf{Локальные LLM:} & Unsloth, LM Studio, Ollama, llama.cpp, OpenRouter \\
    \textbf{Генерация изображений:} & Stable Diffusion, ComfyUI, LoRA, ControlNet \\
    \textbf{Нейросети/ML:} & Перцептроны, CNN, OpenCV, NumPy, Scikit-learn, PyTorch (базовый) \\
    \textbf{Промпт-инжиниринг:} & JSON schema, system prompt, few-shot, chain-of-thought \\
    \textbf{Языки:} & Python, TypeScript, JavaScript, SQL \\
    \textbf{Инструменты:} & Git, Docker, Jupyter Notebook, Google Colab \\
    \textbf{Языки:} & Русский — родной, Английский — A2 (элементарный)
\end{tabular}

% ============================================================
% ДОПОЛНИТЕЛЬНАЯ ИНФОРМАЦИЯ
% ============================================================
\makeFooter{проекты по AI/ML, нейросетям, интеграции LLM}{Интеграция AI-моделей, промпт-инжиниринг, генерация контента, компьютерное зрение, обучение нейросетей}

\end{document}